\documentclass[12pt,a4paper]{article}

% --- Packages ---
\usepackage[utf8]{inputenc}
\usepackage[T1]{fontenc}
\usepackage[french]{babel}
\usepackage{geometry}
\usepackage{enumitem}
\usepackage{array}
\usepackage{fancyhdr}
\usepackage{amsmath, amssymb}

% --- Mise en page ---
\geometry{margin=2.5cm}

% --- Pied de page ---
\pagestyle{fancy}
\fancyhf{}
\fancyfoot[R]{\it{G.M.}}
\renewcommand{\headrulewidth}{0pt}
\renewcommand{\footrulewidth}{0pt}

\setlength{\parindent}{0pt}

\begin{document}

\begin{center}
    \Large \textbf{-- Automatismes - QCM --} \\[5pt]
    \normalsize Aucune justification n’est demandée et une seule réponse est possible 
    par question.
\end{center}

\vspace{0.7cm}

\textbf{Question 1 :}
\vspace{5pt}

Dans une classe de 24 élèves, 6 sont gauchers. La proportion d’élèves gauchers est :
\vspace{0.3cm}

\renewcommand{\arraystretch}{3}
\begin{tabular}{|m{3.4cm}|m{3.4cm}|m{3.4cm}|m{3.4cm}|}
\hline
\centering \vspace{-0.3cm} \textbf{A.} 0,25 & 
\centering \vspace{-0.3cm} \textbf{B.} 
0,4 & 
\centering \vspace{-0.3cm} \textbf{C.} $\dfrac{1}{6}$ & 
\centering \vspace{-0.3cm} 
\textbf{D.} 0,15 \tabularnewline
\hline
\end{tabular}

\vspace{0.7cm}

\textbf{Question 2 :}
\vspace{5pt}

Un article coûte 150 €. Le vendeur propose une réduction de 20\%. Combien d'euros 
représente la réduction ?
\vspace{0.3cm}

\renewcommand{\arraystretch}{3}
\begin{tabular}{|m{3.4cm}|m{3.4cm}|m{3.4cm}|m{3.4cm}|}
\hline
\centering \vspace{-0.3cm} \textbf{A.} 30 € & 
\centering \vspace{-0.3cm} \textbf{B.} 
130 € & 
\centering \vspace{-0.3cm} \textbf{C.} 20 € & 
\centering \vspace{-0.3cm} 
\textbf{D.} 15 € \tabularnewline
\hline
\end{tabular}

\vspace{0.7cm}

\textbf{Question 3 :}
\vspace{5pt}

Dans un collège, 15\% des élèves sont en classe de 3ème, ce qui représente 45 élèves. 
L’effectif total du collège est :
\vspace{0.3cm}

\renewcommand{\arraystretch}{3}
\begin{tabular}{|m{3.4cm}|m{3.4cm}|m{3.4cm}|m{3.4cm}|}
\hline
\centering \vspace{-0.3cm} \textbf{A.} 200 & 
\centering \vspace{-0.3cm} \textbf{B.} 350 & 
\centering \vspace{-0.3cm} \textbf{C.} 400 & 
\centering \vspace{-0.3cm} \textbf{D.} 300 \tabularnewline
\hline
\end{tabular}

\vspace{0.7cm}

\textbf{Question 4 :}
\vspace{5pt}

Un prix augmente de 15\%. Le coefficient multiplicateur associé à cette augmentation 
est :

\vspace{0.3cm}

\renewcommand{\arraystretch}{3}
\begin{tabular}{|m{3.4cm}|m{3.4cm}|m{3.4cm}|m{3.4cm}|}
\hline
\centering \vspace{-0.3cm} \textbf{A.} 0,15 & 
\centering \vspace{-0.3cm} \textbf{B.} 1,15 & 
\centering \vspace{-0.3cm} \textbf{C.} 0,85 & 
\centering \vspace{-0.3cm} \textbf{D.} 1,5 \tabularnewline
\hline
\end{tabular}

\vspace{0.7cm}

\textbf{Question 5 :}
\vspace{5pt}

Un article coûte 80 €. Après une augmentation de 10\%, son nouveau prix est :
\vspace{0.3cm}

\renewcommand{\arraystretch}{3}
\begin{tabular}{|m{3.4cm}|m{3.4cm}|m{3.4cm}|m{3.4cm}|}
\hline
\centering \vspace{-0.3cm} \textbf{A.} 88 € & 
\centering \vspace{-0.3cm} \textbf{B.} 72 € & 
\centering \vspace{-0.3cm} \textbf{C.} 80,1 € & 
\centering \vspace{-0.3cm} \textbf{D.} 90 € \tabularnewline
\hline
\end{tabular}

\vspace{0.7cm}

\textbf{Question 6 :}
\vspace{5pt}

Un article coûte 120 €. Après une réduction de 25\%, son nouveau prix est :
\vspace{0.3cm}

\renewcommand{\arraystretch}{3}
\begin{tabular}{|m{3.4cm}|m{3.4cm}|m{3.4cm}|m{3.4cm}|}
\hline
\centering \vspace{-0.3cm} \textbf{A.} 95 € & 
\centering \vspace{-0.3cm} \textbf{B.} 100 € & 
\centering \vspace{-0.3cm} \textbf{C.} 80 € & 
\centering \vspace{-0.3cm} \textbf{D.} 90 € \tabularnewline
\hline
\end{tabular}

\vspace{0.7cm}

\textbf{Question 7 :}
\vspace{5pt}

Après une augmentation de 20\%, un article coûte 144 €. Son prix initial était :
\vspace{0.3cm}

\renewcommand{\arraystretch}{3}
\begin{tabular}{|m{3.4cm}|m{3.4cm}|m{3.4cm}|m{3.4cm}|}
\hline
\centering \vspace{-0.3cm} \textbf{A.} 110 € & 
\centering \vspace{-0.3cm} \textbf{B.} 120 € & 
\centering \vspace{-0.3cm} \textbf{C.} 130 € & 
\centering \vspace{-0.3cm} \textbf{D.} 100 € \tabularnewline
\hline
\end{tabular}

\vspace{0.7cm}

\textbf{Question 8 :}
\vspace{5pt}

Après une réduction de 30\%, un article coûte 70 €. Son prix initial était :
\vspace{0.3cm}

\renewcommand{\arraystretch}{3}
\begin{tabular}{|m{3.4cm}|m{3.4cm}|m{3.4cm}|m{3.4cm}|}
\hline
\centering \vspace{-0.3cm} \textbf{A.} 110 € & 
\centering \vspace{-0.3cm} \textbf{B.} 90 € & 
\centering \vspace{-0.3cm} \textbf{C.} 100 € & 
\centering \vspace{-0.3cm} \textbf{D.} 80 € \tabularnewline
\hline
\end{tabular}

\vspace{0.7cm}

\textbf{Question 9 :}
\vspace{5pt}

Un prix passe de 50 € à 65 €. Le taux d’évolution, exprimé en pourcentage, est :
\vspace{0.3cm}

\renewcommand{\arraystretch}{3}
\begin{tabular}{|m{3.4cm}|m{3.4cm}|m{3.4cm}|m{3.4cm}|}
\hline
\centering \vspace{-0.3cm} \textbf{A.} 15\% & 
\centering \vspace{-0.3cm} \textbf{B.} 20\% & 
\centering \vspace{-0.3cm} \textbf{C.} 30\% & 
\centering \vspace{-0.3cm} \textbf{D.} 25\% \tabularnewline
\hline
\end{tabular}

\vspace{0.7cm}

\textbf{Question 10 :}
\vspace{5pt}

Un prix subit deux augmentations successives de 10\% puis de 20\%. Le taux d’évolution 
global, exprimé en pourcentage, est :
\vspace{0.3cm}

\renewcommand{\arraystretch}{3}
\begin{tabular}{|m{3.4cm}|m{3.4cm}|m{3.4cm}|m{3.4cm}|}
\hline
\centering \vspace{-0.3cm} \textbf{A.} 30\% & 
\centering \vspace{-0.3cm} \textbf{B.} 32\% & 
\centering \vspace{-0.3cm} \textbf{C.} 27\% & 
\centering \vspace{-0.3cm} \textbf{D.} 35\% \tabularnewline
\hline
\end{tabular}

\end{document}
